\section{Reconstruction from first-moment profiles}
\label{sec:v74-moments}

The separated limiting law has an additional consequence for the amount
of smoothing required by the inverse problem.  Four functions of one
variable determine the two-variable law.  The relevant moment matrix is
nonsingular for a geometric reason: the two marked first derivatives
have opposite signs.  We first prove this assertion on the limiting
model, then retain the finite-flight error, and finally construct an
estimator on the image of actual tables.  All samples in this section
are the same paired endpoint records as before.

\subsection{An exact moment identity}

Fix a symmetric interval $J=[-R_+,R_+]$ and an interior interval
$I=[-R,R]\Subset J$, both containing zero.  For an integrable kernel
$f$ on $J^2$ define
\begin{align}
 M_j(u)&=\int_J v^j f(u,v)\,dv,&
 N_j(v)&=\int_J u^j f(u,v)\,du,\qquad j=0,1,
 \label{eq:v74-profiles}\\
 H_f&=\int_{J^2}\binom{1}{u}(1\ \ v)f(u,v)\,du\,dv.
 \label{eq:v74-matrix}
\end{align}
Write $M(u)=(M_0(u),M_1(u))$ and
$N(v)=(N_0(v),N_1(v))^{\mathsf T}$.  The compression used below is
\[
 \mathcal C_I(f)=\bigl(M_0|_I,M_1|_I,N_0|_I,N_1|_I,H_f\bigr)
 \in\mathcal X_m:=C^m(I)^4\times\mathbb R^{2\times2}.
\]
Equip this space with the sum of the largest function norm and the
matrix operator norm.  The integration domain in all five components
is the full gate $J^2$.  In particular, $H_{f,00}=1$ for a density
normalized on that gate; restriction to $I$ does not renormalize it.
The functions $M_1,N_1$ retain paired information and are not ordinary
marginal densities.

\begin{lemma}[Moment reconstruction]\label{lem:v74-skeleton}
Suppose $f(u,v)=a(u)^{\mathsf T}Cb(v)$, where $a,b$ have two
components and $C$ is a $2\times2$ matrix.  If $H_f$ is nonsingular,
then
\begin{equation}\label{eq:v74-reconstruction}
                  f(u,v)=M(u)H_f^{-1}N(v),\qquad (u,v)\in J^2.
\end{equation}
For each integer $m\ge0$, the map
\[
 \mathcal R:(M_0,M_1,N_0,N_1,H)\longmapsto
                   \bigl[(u,v)\mapsto M(u)H^{-1}N(v)\bigr]
\]
is uniformly Lipschitz into $C^m(I^2)$ on sets with bounded
$\mathcal X_m$ norm and $|\det H|\ge\kappa_*>0$.
Here $C^m(I^2)$ uses derivatives of total order at most $m$.
\end{lemma}
\begin{proof}
Put
\[
 L=\int_J\binom{1}{u}a(u)^{\mathsf T}\,du,
 \qquad R_0=\int_J b(v)(1\ \ v)\,dv.
\]
Fubini's theorem gives $H_f=LCR_0$,
$M(u)=a(u)^{\mathsf T}CR_0$, and $N(v)=LCb(v)$.
Nonsingularity of their square product implies that all three factors
are nonsingular.  Cancelling them proves
\eqref{eq:v74-reconstruction}.  This argument also proves that the
separation rank is exactly two under the hypothesis.

On the stated bounded set, the adjugate formula bounds $\|H^{-1}\|$.
For a second matrix $\widetilde H$ use
\[
 H^{-1}-\widetilde H^{-1}
       =H^{-1}(\widetilde H-H)\widetilde H^{-1}.
\]
Add and subtract the three products obtained by replacing one factor
at a time.  For $i+j\le m$, differentiation of a product gives
$M^{(i)}(u)H^{-1}N^{(j)}(v)$.  The same bounds therefore control
every required derivative.  No derivative of the data above order $m$
is used.
\end{proof}

This is a continuous moment version of the classical exact low-rank
matrix reconstruction underlying CUR decompositions; see
\cite{HammHuang2019}.  Its use here requires, in addition, a
nondegeneracy bound forced by the billiard law and control of the defect
of finite-flight laws from the separated model.

\begin{proposition}[Geometric nondegeneracy]\label{prop:v74-nondegeneracy}
Let the limiting law at one phase be
\begin{equation}\label{eq:v74-law}
 f(u,v)=Z^{-1}w(u)z(v)\{d-A(u)-C(v)\},\qquad (u,v)\in J^2,
\end{equation}
with the hypotheses of \eqref{eq:v72-model}--\eqref{eq:v72-slopes}.
Suppose throughout $J$ that $A'$ has sign $-p$, $C'$ has sign $p$,
and $|A'|,|C'|\ge p_*/2>0$.  Then $\det H_f>0$ and
\eqref{eq:v74-reconstruction} holds.

On a bounded class with $|p|\ge p_*$, fixed $J$, uniform positive
lower and upper bounds for $w,z$, and an upper bound for $Z$, the
determinant has a positive uniform lower bound.  The derivative
hypotheses follow by shrinking $J$ once from the uniform $C^2$ action
bounds and $A'(0)=-p$, $C'(0)=p$.  They are not additional measured
amplitudes or supplied curvature data.
\end{proposition}
\begin{proof}
Use
\[
 a(u)=\binom{w(u)}{w(u)A(u)},\qquad
 b(v)=\binom{z(v)}{z(v)C(v)},\qquad
 C_0=Z^{-1}\begin{pmatrix}d&-1\\-1&0\end{pmatrix}
\]
in Lemma~\ref{lem:v74-skeleton}.  To compute the determinant without
assuming its nonvanishing, set $W=\int_Jw$, $V=\int_Jz$, and let
$\mu_w,\mu_z$ have densities $w/W,z/V$.  Direct evaluation gives
\begin{equation}\label{eq:v74-determinant}
 \det H_f=-\frac{W^2V^2}{Z^2}
  \operatorname{Cov}_{\mu_w}(u,A(u))
  \operatorname{Cov}_{\mu_z}(v,C(v)).
\end{equation}
Indeed the two covariance factors, multiplied by $W^2,V^2$, are
respectively $\det L,\det R_0$, and $\det C_0=-Z^{-2}$.
For a probability measure $\mu$ and independent $X,Y$ of law $\mu$,
\[
 \operatorname{Cov}_{\mu}(x,g(x))
       =\tfrac12\mathbb E[(X-Y)(g(X)-g(Y))].
\]
Consequently a derivative of constant sign and magnitude at least
$p_*/2$ gives a covariance of that sign and magnitude at least
$(p_*/2)\operatorname{Var}_{\mu}(x)$.  The two covariances in
\eqref{eq:v74-determinant} have opposite signs.

For completeness, if $w_*\le w\le w^*$ on $J$, then
\[
 \operatorname{Var}_{\mu_w}(u)
 \ge\frac{w_*}{2R_+w^*}\inf_{c\in\mathbb R}
                      \int_{-R_+}^{R_+}(u-c)^2\,du
 =\frac{w_*}{w^*}\frac{R_+^2}{3}.
\]
The same estimate applies to $z$, while $W,V$ have positive lower
bounds.  Substitution in \eqref{eq:v74-determinant} proves the uniform
claim.  Finally, if the second action derivatives are bounded by $K_2$,
choose $K_2R_+\le p_*/2$.  The mean value theorem gives the asserted
sign and magnitude bounds.  All constants may deteriorate as the fixed
gate shrinks or $p_*$ tends to zero; no uniformity in those limits is
claimed.
\end{proof}

\begin{corollary}[Moment determination of the local invariant]
\label{cor:v74-invariant}
For the oblique marked class, four limiting moment profiles and their
matrix $H_f$ at each phase determine the offsets and all smooth contact
germs.  Adding the charged success probabilities as a function of the
returning length determines the free area.  Conversely, common contact
germs, offsets and free area give identical compressed charged local
experiments on sufficiently small common gates.
\end{corollary}
\begin{proof}
Choose the common gates as in Proposition~\ref{prop:v74-nondegeneracy}.
The moment identity reconstructs each limiting law on $I^2$.
Theorem~\ref{thm:v72-rigidity} determines the offsets and smooth
contact germs.  Their curvatures give $D_{N,b}$, and
\eqref{eq:v73-count-scalar} then determines the area from one phase.
The converse is the locality direction of
Theorem~\ref{thm:v73-fibers}, followed by integration.  Applying this
argument on a cofinal collection of gates gives the assertion for germs.
\end{proof}

At finite $N$ the identity in Lemma~\ref{lem:v74-skeleton} need not
hold for $f_N$.  Thus the corollary concerns a family of laws and its
relative limit, not a sufficient statistic for one finite sample.
Nor does it replace paired endpoints by two unpaired marginal samples.
For example, at normal incidence the positive model
$Z^{-1}(d-a u^2-c v^2)$ on a symmetric square can have rank two but
$M_1=N_1=0$, so $H_f$ is singular.  The two-offset results at normal
incidence remain available independently of the first-moment inverse.

\subsection{Finite-flight comparison on the physical image}

Fix the bounded class $\mathscr U$ of
Proposition~\ref{prop:v73-finite-record-stability}.  In particular,
$m\ge3$ satisfies the explicit later-visit condition on the enlarged
positive quadratic class.  Choose $J$, $I\Subset J$, and a contact
collar $I_c\Subset I$ once so that the preceding proposition and the
inherited action-to-contact stability apply uniformly.  Densities are
normalized on the success gate $J^2$; all integrations below still use
that whole gate.  The relative physical theorem supplies
\begin{equation}\label{eq:v74-relative}
 \max_b\|f_{N,b}^{\vartheta}-f_b^{\vartheta}\|_{C^m(J^2)}
                     \le C_0e^{-\omega N},\qquad N\ge N_0,
\end{equation}
uniformly over actual pairs $\vartheta=(T,d)\in\mathscr U$.
Shrinking the fixed gates changes class constants, not the charged
preparation convention.

\begin{proposition}[Stability of compressed finite records]
\label{prop:v74-finite-stability}
For two actual pairs in $\mathscr U$ put
\[
 \Delta_M=\max_b\|\mathcal C_I(f_{N,b}^{\vartheta})-
                   \mathcal C_I(f_{N,b}^{\widetilde\vartheta})\|_{\mathcal X_m},
 \quad e_N=\Delta_M+e^{-\omega N},
 \quad\Delta_\pi=\left|\log
     \frac{\pi_{N,b_0}^{\widetilde\vartheta}}{\pi_{N,b_0}^{\vartheta}}\right|.
\]
There are class constants $C,\eta_*>0$ such that $e_N\le\eta_*$
implies
\begin{align}
 \|F^T-F^{\widetilde T}\|_{C^0(I_c)}
  +\|d-\widetilde d\|_\infty+\|c-\widetilde c\|_\infty
                &\le Ce_N,\label{eq:v74-profile-stability}\\
 \left|\log\frac{\mathcal A(\widetilde T)}{\mathcal A(T)}\right|
                &\le C(1+N)e_N+\Delta_\pi.
                \label{eq:v74-area-stability}
\end{align}
Here $c$ is the vector of contact curvatures.  The estimate does not
assume exact rank two at finite flight length.
\end{proposition}
\begin{proof}
Integration in \eqref{eq:v74-profiles}--\eqref{eq:v74-matrix} is a
bounded linear map from $C^m(J^2)$ to $\mathcal X_m$.  Hence the two
limiting compressed tuples are at distance at most
$\Delta_M+Ce^{-\omega N}$.  Proposition~\ref{prop:v74-nondegeneracy}
gives their uniform determinant floor.  Lemma~\ref{lem:v74-skeleton}
therefore bounds the distance of the two limiting densities on $I^2$
by $Ce_N$.  Equation~\eqref{eq:v74-relative} also gives
\begin{equation}\label{eq:v74-finite-density-control}
 \max_b\|f_{N,b}^{\vartheta}-f_{N,b}^{\widetilde\vartheta}
                                             \|_{C^m(I^2)}\le Ce_N.
\end{equation}
Now apply Proposition~\ref{prop:v73-finite-record-stability} on
these fixed interior and contact collars.  Its proof uses the
$C^m$ law norm to recover the actions and their quadratic coefficients
before estimating the profiles.  In particular it does not infer
curvature control from $C^0$ profile control.  The exact same-gate
central-density identity and the logarithmic twist comparison then give
\eqref{eq:v74-area-stability}.

Equivalently, the preceding argument shows that
$\mathcal R(\mathcal C_I(f_{N,b}))$ approximates $f_{N,b}|_{I^2}$
with error $Ce^{-\omega N}$, for $N$ large enough.  This is only an
approximation: the finite central density in the area identity is the
one belonging to the actual candidate, not an exactly identified
rank-two surrogate.
\end{proof}

\subsection{Estimation of the four profiles}

We give the sampling estimate in detail because smoothing a weighted
marginal, rather than a two-variable density, changes the exponents.
Let $s\ge1$ be an integer.  Assume the same uniform forward bounds
as before, now stated on the full fixed gate:
\[
 \sup_{N\ge N_0,\,b,\,\vartheta\in\mathscr U}
                   \|f_{N,b}^{\vartheta}\|_{C^{m+s}(J^2)}\le K.
\]
All constants below may depend on this class, the fixed intervals,
$m,s$ and the number $r$ of phases.

Choose $K_0\in C_c^{m+1}((-1,1))$ with integral one and vanishing
moments of orders $1,\ldots,s-1$.  Such a kernel can be constructed
by multiplying a positive smooth bump by a polynomial of degree at
most $s-1$: the moment conditions are a nonsingular Gram system for
polynomials on the support of that bump.  The kernel need not be
nonnegative.  Let $h_0\le1$ be smaller than the distance from $I$ to
$\partial J$.

For $n$ accepted pairs at one phase, write their true coordinates as
$(U_i,V_i)$ and their readouts as $(\widetilde U_i,\widetilde V_i)$,
with both coordinate errors at most $0\le\delta\le1$.  Define
\begin{align}
 \widehat M_j(u)&=\frac1{nh}\sum_{i=1}^n\widetilde V_i^j
                   K_0((u-\widetilde U_i)/h),\label{eq:v74-est-M}\\
 \widehat N_j(v)&=\frac1{nh}\sum_{i=1}^n\widetilde U_i^j
                   K_0((v-\widetilde V_i)/h),\qquad j=0,1,\\
 \widehat H&=\frac1n\sum_{i=1}^n
                      \binom{1}{\widetilde U_i}(1\ \ \widetilde V_i).
                      \label{eq:v74-est-H}
\end{align}
Denote the resulting tuple by $\widehat{\mathcal C}_{N,b}$.  There
is no inversion of a random nearly singular matrix in this definition.

\begin{lemma}[One-dimensional smoothing bound]
\label{lem:v74-estimation}
For $n\ge2$, $n^{-1}\le h\le h_0$, and $0<\zeta<1/2$, put
$L_n=\log(C_1rn/\zeta)$.  The first $n$ accepted latent pairs at
each phase may be used in \eqref{eq:v74-est-M}--\eqref{eq:v74-est-H}.
With probability at least $1-\zeta$,
\begin{equation}\label{eq:v74-kernel-bound}
 \begin{split}
 \max_b\|\widehat{\mathcal C}_{N,b}-
                 \mathcal C_I(f_{N,b})\|_{\mathcal X_m}
 \le C\left\{h^s+\sqrt{\frac{L_n}{nh^{2m+1}}}
       +\frac{L_n}{nh^{m+1}}+n^{-2}
       +\delta h^{-(m+2)}\right\}.
 \end{split}
\end{equation}
The perturbations may depend arbitrarily on the latent sample.  The
probability statement is about the independent latent preparations;
acceptance and its success/failure tags are exact.
\end{lemma}
\begin{proof}
First set $\delta=0$.  For $0\le k\le m$, the expectation of
$\widehat M_j^{(k)}(u)$ is the convolution of $M_j^{(k)}$ with the
rescaled kernel on the interior interval.  Integration by parts is
legitimate because the kernel support does not meet $\partial J$.
The $C^{m+s}$ bound on $f_N$ gives the same order bound on $M_j$.
Taylor's formula and the moment cancellations give a uniform bias
$Ch^s$.  The argument for $N_j$ is identical.

For a fixed $u$, a summand before division by $n$ has magnitude at
most $Ch^{-k-1}$ and variance at most $Ch^{-2k-1}$.  Indeed the
integral of the square of $K_0^{(k)}((u-U)/h)$ against the bounded
joint density gains one factor $h$; $|V|^j$ is uniformly bounded on
$J$.  Bernstein's inequality thus gives, with failure probability
at most $2e^{-t}$, fluctuation at most
\[
 C\left\{\sqrt{\frac{t}{nh^{2k+1}}}
                    +\frac{t}{nh^{k+1}}\right\}.
\]
Take a grid in $I$ of mesh at most $n^{-2}h^{m+2}$.  Both the
empirical derivative and its expectation have Lipschitz constants
at most $Ch^{-m-2}$.  Passage from the grid to the whole interval
therefore costs at most $Cn^{-2}$.  Since $h\ge n^{-1}$, the grid
has at most $C n^{m+4}$ points.  A union bound over the four profiles,
$r$ phases and $m+1$ derivative orders, with $t=C_m L_n$, proves
the stated stochastic bound.  The entries of $\widehat H$ are
averages of bounded variables and have error at most
$C\{\sqrt{L_n/n}+L_n/n\}$ on the same enlarged event.  This is
absorbed by the displayed bound when $h\le1$.  Independence of the
four profiles is neither assumed nor needed.

Restore the readout errors.  The mean value theorem for the kernel
and the bounded weights gives, for each summand of a derivative of
order $k$,
\[
 \left|\widetilde V_i^j h^{-k-1}
 K_0^{(k)}((u-\widetilde U_i)/h)
 -V_i^j h^{-k-1}K_0^{(k)}((u-U_i)/h)\right|
 \le C\delta\{h^{-k-2}+h^{-k-1}\}.
\]
The weights remain bounded by $C(R_+)$ because $\delta\le1$.
For $k\le m$ and $h\le1$ the last bound is
$C\delta h^{-m-2}$.  Averaging does not increase it.  The matrix
entries have deterministic perturbation at most $C\delta$, also
absorbed.  This comparison holds sample by sample and so requires
no distributional assumption on the readout errors.
\end{proof}

\subsection{Charged recovery with one-dimensional smoothing}

The actual finite image to be fitted is
\begin{equation}\label{eq:v74-image}
 \mathcal K_N(\vartheta)=
  \bigl((\mathcal C_I(f_{N,b}^{\vartheta}))_{b=0}^{r-1},
                             \log\pi_{N,b_0}^{\vartheta}\bigr)
                    \in\mathcal X_m^r\times\mathbb R.
\end{equation}
Use the maximum of the tuple norms plus the scalar absolute value.
For each returning $N$, fix a countable dense subset of this image,
and one actual representative $(T_j,d_j)\in\mathscr U$ for each
of its points.  This is possible because $C^m(I)$, and hence the
ambient metric space, is separable.  Closedness of the image is not
required.

Perform $B$ independent normalized unit-speed Liouville resets at each
phase.  Put $q_N=c_0e^{-\Gamma N}$ and $n=\lfloor Bq_N/2\rfloor$.
If every success count $S_b$ is at least $n\ge2$, use the first $n$
accepted pairs for the estimated compressed tuples and set
$\widehat\pi=S_{b_0}/B$.  Select the first representative whose
distance to
$((\widehat{\mathcal C}_{N,b})_b,\log\widehat\pi)$ is within
$n^{-2}$ of the infimum over the countable image.  Otherwise select a
fixed member of $\mathscr U$.  The infimum of countably many
measurable distances is measurable, and a qualifying index exists by
the definition of infimum.  Thus this specifies an actual table-offset
estimator on every record.  Its area is always the area of that same
table.  An effective enumeration or polynomial-time search is not
asserted.

\begin{theorem}[Compressed charged reconstruction]
\label{thm:v74-charged}
In the stated class set
\begin{equation}\label{eq:v74-exponents}
 \alpha_1=\frac{s}{2(m+s)+1},\qquad
 \gamma_1=\frac{s}{m+s+2}.
\end{equation}
Suppose $N\ge N_0$, $n\ge2$, $Bq_N\ge C_2\log(4r/\zeta)$,
and $L_n/n\le1$.  Choose
\[
 h=\max\{(L_n/n)^{1/[2(m+s)+1]},
                      \delta^{1/(m+s+2)},n^{-1}\}\le h_0.
\]
For $\varepsilon_N=e^{-\omega N}+(L_n/n)^{\alpha_1}
+\delta^{\gamma_1}$ sufficiently small, the estimator just defined
satisfies, with probability at least $1-2\zeta$,
\begin{align}
 \|F^{\widehat T}-F^T\|_{C^0(I_c)}
 +\|\widehat d-d\|_\infty+\|\widehat c-c\|_\infty
                 &\le C\varepsilon_N,\label{eq:v74-charged-local}\\
 \left|\log\frac{\mathcal A(\widehat T)}{\mathcal A(T)}\right|
                 &\le C(1+N)\varepsilon_N.
                 \label{eq:v74-charged-area}
\end{align}
All $rB$ preparations, including failures, are charged.  No additional
phase group is needed for the area.  Readout errors occur after exact
acceptance and do not corrupt the success tags.
\end{theorem}
\begin{proof}
The count argument in \eqref{eq:v73-log-count}, with
$t=\log(4r/\zeta)$, gives simultaneously
\[
 S_b\ge n\quad\hbox{for every }b,\qquad
 |\log\widehat\pi-\log\pi_{N,b_0}|
                     \le C\sqrt{t/(Bq_N)}
\]
except on an event of probability at most $\zeta$.  The
Bernoulli/conditional-mark representation of independent resets gives
independent latent marks of density $f_{N,b}$ at successive successes.
They may be extended beyond the finite record when proving a probability
bound; on the count event only the first $n$ actually observed marks
are used.  This construction does not condition the marks on a
possibly mark-dependent noisy acceptance test: acceptance is exact.

Apply Lemma~\ref{lem:v74-estimation} to these latent sequences.
At the prescribed bandwidth, its right side is at most
$C\{(L_n/n)^{\alpha_1}+\delta^{\gamma_1}\}$.
For the variance term this follows by substituting the first bandwidth
lower bound; the envelope exponent becomes
$(m+2s)/[2(m+s)+1]\ge\alpha_1$.  The readout term is bounded using
the second lower bound, and the bias is bounded using the maximum
formula.  The $n^{-2}$ term is absorbed.  Since $\alpha_1<1/2$,
$n$ is comparable to $Bq_N$, and $L_n/n\le1$, the count error is
absorbed by the same expression.

Density of the countable actual image implies that its infimal
distance to the estimated tuple is no larger than the distance of the
true image.  The selected image is therefore within twice the
estimation error plus $n^{-2}$ of the true image.  Both candidates
are actual pairs.  Proposition~\ref{prop:v74-finite-stability}
gives the assertions, with the physical error $e^{-\omega N}$ added
only at this step.  The union of the count and smoothing exceptional
events has probability at most $2\zeta$.
\end{proof}

\begin{corollary}[Budget-dependent flight length]
\label{cor:v74-budget}
Let
\[
 \beta_1=\frac{\alpha_1\omega}{\omega+\alpha_1\Gamma},\qquad
 L_B=\log(C_*rB/\zeta),\qquad
 t=(L_B/B)^{\beta_1}+\delta^{\gamma_1}<1,
\]
and choose $N=P\lfloor\log(1/t)/(P\omega)\rfloor$.
Whenever the hypotheses of Theorem~\ref{thm:v74-charged} hold,
its local error is at most $Ct$ and its joint local-and-area error
is at most
\begin{equation}\label{eq:v74-budget}
                            Ct\log(e/t).
\end{equation}
\end{corollary}
\begin{proof}
The floor gives $e^{-\omega N}\le e^{\omega P}t$,
$1+N\le C\log(e/t)$ and $n\ge cBt^{\Gamma/\omega}$ under the
count conditions.  Write $x=(L_B/B)^{\beta_1}\le t$.  Since
$L_n\le CL_B$,
\[
 (L_n/n)^{\alpha_1}
 \le C(L_B/B)^{\alpha_1}t^{-\alpha_1\Gamma/\omega}
 \le C(L_B/B)^{\alpha_1}x^{-\alpha_1\Gamma/\omega}=Cx.
\]
The last equality is the definition of $\beta_1$.
Also $\delta^{\gamma_1}\le t$.  Substitution in
\eqref{eq:v74-charged-local}--\eqref{eq:v74-charged-area} proves
the claims.
\end{proof}

For the same $m,s,\omega,\Gamma$, the exponents $\alpha_1$ and
$\gamma_1$ strictly improve the two-dimensional smoothing exponents
$s/[2(m+s)+2]$ and $s/(m+s+3)$ in
Theorem~\ref{thm:v73-charged-invariant}.  The corresponding budget
exponent improves as well, since
$\alpha\mapsto\alpha\omega/(\omega+\alpha\Gamma)$ is strictly
increasing.  These are comparisons of proved upper bounds on the same
oblique bounded class, not minimax assertions or comparisons with lens
or marked-length observations.  The improvement comes from estimating
one-dimensional weighted profiles, not from dividing a small absolute
error by a rare success probability.  The full smooth inverse and its
finite-order later-visit estimate remain essential to the passage from
these profiles to the actual reflecting contacts.
